\topic{Sources, authorship, and review limits}
\labelled{Established mathematics; newly authored activity.}{All prose, examples, diagrams and solution arguments in this guide were newly composed for Week 64 around established mathematics. No source page, source figure, borrowed problem, or borrowed writing example is reproduced. Path starts, comparisons and facilitation are activity design choices; the references do not claim classroom testing of these materials.}
\begin{enumerate}[label=\textbf{S\arabic*.},leftmargin=25pt]
\item Alex Wright, \emph{From Rational Billiards to Dynamics on Moduli Spaces}, pp. 3--6: translation gluing, the regular octagon, its single \(6\pi\) cone point and genus 2. P. 10 gives advanced dynamics, not a conclusion from tracing.\par
\url{https://public.websites.umich.edu/~alexmw/BilliardsToModuli.pdf}
\item The surface-dynamics authors, \emph{Square-tiled Surfaces}, official documentation, First steps and Topological invariants. The three-square example with permutations \((1,2)\), \((1,3)\) verifies the A/B right swap, A/C up swap and genus-two one-singularity model.\par
\url{https://flatsurf.github.io/surface-dynamics/examples/square_tiled_surfaces.html}
\item Giovanni Forni and Carlos Matheus, \emph{Introduction to Teichm\"uller Theory and Its Applications to Dynamics of Interval Exchange Transformations, Flows on Surfaces and Billiards}, pp. 17--18, Figure 4: the three-square L and translation structure; p. 78: horizontal cylinders of widths 1 and 2. Our exact starts and paths are independently calculated.\par
\url{https://www.math.uchicago.edu/~masur/gm.pdf}
\item Jeff Erickson, \emph{Surface Homotopy}, Universal Cover: the hyperbolic genus-two octagon picture. Used to distinguish that metric from the Euclidean cone surface.\par
\url{https://jeffe.cs.illinois.edu/teaching/comptop/2020/notes/21-surface-homotopy.html#universal-cover}
\end{enumerate}
\labelled{Verification status.}{Prepared October 5, 2026. Translations, specified trajectories, first returns, corner classes and Euler counts are checked mathematically. Final digital review and student-page alignment are recorded with the editable source. This is not classroom piloting or physical rehearsal.}
\topic{After a session: preserve useful evidence}
Record which problems were used, not just printed. Keep one checked example and one moment of difficulty. Note whether children controlled the next segment, matched points unaided, kept headings parallel and recognized a complete return. Separate what happened from explanations for why. On a return visit, reuse the familiar rule with a new question rather than requiring unfinished pages to be completed in order.
